\documentclass{article}
\usepackage[utf8]{inputenc}
\usepackage[T1]{fontenc}
\usepackage[margin=2cm]{geometry}
\usepackage{hyperref}
\sloppy
\begin{document}
\section*{Historical source}
Encyclopaedia Britannica scan, source row 2473. Parse 3.5 extracted the text. Super 3.5 VL corrected it from an enlarged crop of each block. It can still contain recognition errors.
Page headers, footers, and page numbers are removed.

\subsection*{Untitled continuation (E0231)}
\textit{Continues an entry from an earlier page.}

ERA (sit)

The onager, or wild ass, has, by some authors, been confounded with the zebra; but very improperly, for this last is a distinct species; for the onager is not freaked like this, nor is his shape so beautiful. Wild asses are said to be very swift of course; and when they see a man, they make a bound, and immediately fly away; insomuch, that there is no taking of them, but by traps and gins. They have much the same shape as common asses; but they are of a brighter colour, and there runs a white lift from the head to the tail. Of the hide of these asses, and particularly of that part next the rump, they make that excellent leather which we call shagreen, and which is put to so many curious uses.

In America they have no asses at all, nor yet horses; but they have been carried thither long ago, at first by the Spaniards, and afterwards by other nations, where they have multiplied greatly; insomuch, that, in some places, there are whole droves of them that run wild, and are very hard to be caught. Asses in general carry the heaviest burdens in proportion to their bulk; and, as their keeping costs little or nothing, it is a great wonder they are not put to more uses than they generally are among us.

The flesh of the common ass is never eaten in these parts of the world; though some pretend their colts are tender, and not disagreeable. [Plate LXXV. fig. 2.]

3. The zebra.—This animal has the figure and gracefulness of the horse, joined to the swiftness of the flag. He is about 7 feet long, from the point of the muzzle to the origin of the tail, and about 4 feet high. The colour of his skin is beautiful and uniform, confifting of alternate parallel rings of black and white disposed in the most regular manner, as reprefented in the plate, [LXXV. fig. 3.] He is generally less than the horse, and larger than the ass.

the Zebra is found no where but in the eastern and southern provinces of Africa, from Ethiopia to the Cape of Good Hope, and from the Cape of Good Hope to Congo. The Dutch have been at great pains to tame and use them for domestic purposes, but with little success. He is hard-mouthed, and kicks when any person attempts to touch or come near him. He is restless and obfuscate as a mule: but perhaps the wild horse is naturally as untractable as the Zebra; for, it is probable, if he were early accustomed to obedience and a domestic life, he would become as docile as the horse.

\subsection*{ERANARCHIA (E0232)}
ERANARCHIA, a public officer among the ancient Greeks, whose business was to preside over and direct the alms and provisions made for the poor. Cornelius Nepos, in his life of Epaminondas, describes his office thus: when any person was reduced to poverty, taken captive, or had a daughter to marry, which he could not effect for want of money, \&c. the eranarcha called an assembly of friends and neighbours, and taxed each according to his means and estate, to contribute towards his relief.

\subsection*{ERANTHEMUM (E0233)}
ERANTHEMUM, in botany, a genus of the diandria monogynia class. The calix is divided into five segments; the tube is filiform; and the fligma is limple. There is but one species, a native of Ethiopia.

\subsection*{ERASED (E0234)}
ERASED, in heraldry. See ARRACHE.

\subsection*{ERECTION (E0235)}
ERECTION, in a general sense, the art of raising or e-
levating any thing, as the erection of a perpendicular,
\&c. It is also used in a figurative sense, as the erec-
tion of a bishopric, marquisate, \&c.

ERECTOR clitoris. See Clitoris.

\subsection*{ERECTOR PENIS (E0236)}
ERECTOR PENIS. See PENIS.
EREMIT. See HERMIT.

\subsection*{EREMIT (E0237)}
EREMIT. See HERMIT.

\subsection*{EREMITA (E0238)}
EREMITA, in zoology. See SCARABAEUS.

\subsection*{ERFURT (E0239)}
ERFURT, a large and beautiful city of Upper Saxony in Germany, capital of Thuringia, and subject to the elector of Mentz: E. long. 11° 6', N. lat. 51°.

\subsection*{ERGOT (E0240)}
ERGOT, in farriery, is a stub, like a piece of soft horn, about the bigness of a chestnut, placed behind and below the pastern-joint, and commonly hid under the tuft of the fetlock.

\subsection*{ERICA (E0241)}
ERICA, or Heath, in botany, a genus of the octan-
dria monogynia class. The calix consists of four
leaves, and the corolla of four segments; the filaments
are inserted into the receptacle; the anthers are bifid;
and the capsule has four cells. There are thirty-eight
species, five of which are natives of Britain, viz. the
vulgaris, or common heath; the cinerea, or fine-leaved
heath; the tetralix, or cross-leaved heath; the ciliaris;
or rough-leaved heath; and the multiflora, or fire-
leaved heath.

\subsection*{ERIDANUS (E0242)}
ERIDANUS, in astronomy. See Vol. I. p. 487.

\subsection*{ERIE (E0243)}
ERIE, a salt lake to the westward of Pennsylvania, in North America, situated between 80° and 87° W. long. and between 41° and 42° N lat.

\subsection*{ERIGERON (E0244)}
ERIGERON, or SWEET FLEA BANE, in botany, a genus of the syngenesia polygamia superflua class. The receptacle is naked; the pappus is hairy; and the radii of the corolla are linear and very narrow. There are fifteen species, two of which are natives of Britain, viz. the acre, or blue-flowered flea-bane; and the canadense, or Canada flea-bane.

\subsection*{ERINACEUS (E0245)}
ERINACEUS, or Hedge-hog, in zoology, a genus of quadrupeds belonging to the order of seræ, the characters of which are these: they have two foreteeth in the upper jaw, at a considerable distance from one another, and two in the under jaw, left distant; and they have two recumbent dog-teeth, one on each side. There are three species, viz. 1. The europæus, or common hedge-hog, with round ears, and creted no-frils. It is about nine inches long; the upper part of the body is totally covered with sharp prickles, and the under part is covered with hair. The hedge-hog, even when flanding on his legs, has a very ugly aspect. His body is an oblong mass, convex above, terminated on the fore-part by a very sharp muzzle, and mounted on four short legs, of which nothing appears but the feet, and the tail is not discernible. His ears are broad, round, and short; and his eyes are small and protuberant. The length of his body, from the point of the muzzle to the anus, is about nine inches.

The hedge-hog has a very uncommon method of defending himself from the attacks of other animals: being possessed of little strength or agility, he does not attempt to fly from or assail his enemies; but erects his bristles, and rolls himself up like a ball, exposing no-


\section*{Source record}
\begin{description}
\item[Scan ID] \texttt{britannica:309f765b1ecc83e0a7f7762e44b4a927a7555443:train:2473}
\item[Image file] \texttt{page-2473-b849fb7f25f2.jpg}
\item[Image SHA-256] \texttt{b849fb7f25f25b39685069d12c13b17d47169beeea958446bdff439c230a0f65}
\item[Dataset] NationalLibraryOfScotland/encyclopaedia\_britannica\_illustrated at revision \texttt{309f765b1ecc83e0a7f7762e44b4a927a7555443} (CC0-1.0)
\item[Extraction] model \texttt{nvidia/nemotron-3.5-parse}, raw reply \texttt{extract/raw/page-2473-b849fb7f25f2.json}
\item[Raw reply SHA-256] \texttt{55870fc531bf0bfb53e7079c7244a83080c4a49ec47639f96965a611b24f652d}
\item[Correction] model \texttt{nvidia/Nemotron-3.5-Super-VL}, block crops, raw replies \texttt{correct/raw/crop\_nothink/page-2473-b849fb7f25f2-b*.json}; super\_vl: 22
\item[Running header] none
\end{description}
\end{document}
